
\textbf{Lemme 1.7.} \uline{Soient} $A$ \uline{un anneau principal}, $X=\Spec(A)$. \uline{Alors l'homomorphisme canonique} :
\[
\mathbb Z\longrightarrow K_\bullet(X)
\]
\uline{est un isomorphisme}.

Comme $X$ est régulier, l'homomorphisme $K^{\bullet}(X)\to K_\bullet(X)$ est un isomorphisme. Donc $K_\bullet(X)$ est engendré par les classes des $\underline{O}_X$-Modules localement libres de type fini, i.e.\ par les classes des $A$-modules projectifs de type fini. Comme $A$ est principal, tout module projectif de type fini est libre, d'où la surjectivité. L'injectivité vient de ce que $K^{\bullet}(X)$ est $\mathbb Z$-augmenté.

\textbf{Proposition 1.8.} \uline{Soient} $S$ \uline{un schéma noethérien}, $\underline{E}$ \uline{un} $\underline{O}_S$\uline{-Module localement libre de rang} $r$, $X$ \uline{le fibré vectoriel associé à} $\underline{E}$, $f$ \uline{le morphisme structural de} $X$, \uline{et} $g$ \uline{la section nulle de} $X$. \uline{Alors l'homomorphisme composé} :
\[
K_\bullet(S)\xrightarrow{g_*}K_\bullet(X)\xrightarrow{g^*}K_\bullet(S)
\]
(\uline{où} $g^*$ \uline{est défini d'après} IV 2.12.1 \uline{et} 1.5) \uline{est la multiplication par l'élément} $\lambda_{-1}(E)$ \uline{de} $K^{\bullet}(S)$, \uline{en notant} $E$ \uline{la classe de} $\underline{E}$ \uline{dans} $K^{\bullet}(S)$, \uline{et en posant}
\[
\lambda_{-1}(E)=\sum_{i=0}^r(-1)^i\lambda^i(E).
\]

Cet énoncé est l'analogue, pour le $K_\bullet$, de VII 2.7.

Soit $\underline{F}$ un $\underline{O}_S$-Module cohérent, de classe $F$ dans $K_\bullet(S)$. Alors $g_*(F)$ est la classe dans $K_\bullet(X)$ du complexe $\mathbb Rg_*(\underline{F})=g_*(\underline{F})$, donc de $\underline{F}$ considéré comme $\underline{O}_X$-Module par l'intermédiaire de $g$. Pour calculer $g^*(g_*(f))$, il faut donc prendre la classe de \(\mathbb Lg^*(\underline{F})=\underline{F}\otimes^{\mathbb L}_{\underline{O}_X}\underline{O}_S\).

